\DocumentMetadata{
  pdfversion=2.0,
  pdfstandard=ua-2,
  lang=en-US,
  tagging=on,
  tagging-setup={math/setup=mathml-SE,tabsorder=structure}
}
\documentclass[11pt,letterpaper]{article}
\usepackage[margin=0.68in,includefoot]{geometry}
\usepackage{newcomputermodern}
\usepackage{amsmath,amscd,mathtools}
\usepackage{tikz,adjustbox}
\providecommand{\square}{\mdlgwhtsquare}
\usepackage[hidelinks]{hyperref}
\hypersetup{
  pdftitle={Probability PhD Exam, September 2009},
  pdfauthor={University of Florida Department of Mathematics},
  pdfsubject={Graduate mathematics examination},
  pdfdisplaydoctitle=true,
  bookmarksopen=true
}
\setcounter{secnumdepth}{0}
\setlength{\parindent}{0pt}
\setlength{\parskip}{0.28em}
\setlength{\emergencystretch}{3em}
\setlength{\leftmargini}{2.15em}
\setlength{\leftmarginii}{2.65em}
\setlength{\leftmarginiii}{3em}
\pagestyle{empty}
\raggedbottom
\allowdisplaybreaks
\NewTaggingSocketPlug{block/list/label}{examproblem}{%
  \tagstructbegin{tag=Lbl}\tagstructend%
  \tagstructbegin{tag=\LBody}%
  \tagstructbegin{tag=\ProblemHeadingTag,title-o={\ProblemHeadingText}}%
  \tagmcbegin{tag=\ProblemHeadingTag,actualtext={\ProblemHeadingText}}%
  #2%
  \tagmcend%
  \tagstructend%
}
\newenvironment{examproblems}[2]{%
  \def\ProblemHeadingTag{#2}%
  \AssignTaggingSocketPlug{block/list/label}{examproblem}%
  \begin{enumerate}%
  \setcounter{enumi}{#1}%
  \renewcommand{\labelenumi}{\textbf{\theenumi.}}%
  \setlength{\topsep}{0.35em}%
  \setlength{\partopsep}{0pt}%
  \setlength{\itemsep}{0.48em}%
  \setlength{\parsep}{0.18em}%
}{\end{enumerate}}
\newenvironment{examsubparts}{%
  \AssignTaggingSocketPlug{block/list/label}{default}%
  \begin{enumerate}%
  \setlength{\topsep}{0.18em}%
  \setlength{\partopsep}{0pt}%
  \setlength{\itemsep}{0.16em}%
  \setlength{\parsep}{0.1em}%
}{\end{enumerate}}
\newenvironment{examinstructions}{%
  \par\begingroup\itshape
}{%
  \par\endgroup\vspace{0.18em}
}
\makeatletter
\renewcommand\subsection{\@startsection{subsection}{2}{0pt}%
  {-1.25ex plus -.25ex minus -.1ex}%
  {0.55ex plus .1ex}%
  {\normalfont\large\bfseries}}
\renewcommand{\@maketitle}{%
  \newpage\null\vskip -1em%
  {\footnotesize\bfseries
    \href{https://www.ufl.edu/}{University of Florida}, \href{https://math.ufl.edu/}{Department of Mathematics}\par}%
  \vskip -0.5em%
  \tagstructbegin{tag=H1,title={Probability PhD Exam, September 2009}}%
  \tagpdfparaOff%
  \tagmcbegin{tag=H1}%
    {\LARGE\bfseries\href{https://gma.math.ufl.edu/exams/probability-phd/}{Probability PhD Exam}, September 2009\par}%
  \tagmcend%
  \tagpdfparaOn%
  \tagstructend%
  \vskip 0.25em%
  \tagmcbegin{artifact}%
    \hrule height 1pt%
  \tagmcend%
  \vskip 1em%
}
\makeatother
\title{Probability PhD Exam, September 2009}
\author{}
\date{}
\begin{document}
\maketitle
\thispagestyle{empty}
\begin{examproblems}{0}{H2}
\def\ProblemHeadingText{Problem 1.}
\item
For each fixed~\(n\), \(X_n^1, X_n^2, \ldots, X_n^n\) are independent random variables and have the same Binomial distribution, that is, \(\mathbb P(X_n^i=1)=p_n\) and \(\mathbb P(X_n^i=0)=1-p_n\) for \(i=1,2,\ldots,n\). Let \(S_n=X_n^1+X_n^2+\cdots+X_n^n\). Given the condition that \(np_n\) converges to a positive constant~\(\beta\) as~\(n\) goes to~\(\infty\), find the limit distribution of \(S_n\) and prove it.
\def\ProblemHeadingText{Problem 2.}
\item
Let~\(X\) and~\(Y\) be two independent random variables with \(\mathbb E[Y]=0\). Show that for \(p\geq 1\), we have
\[
\mathbb E|X+Y|^p\geq \mathbb E|X|^p.
\]
\def\ProblemHeadingText{Problem 3.}
\item
Let~\(X\) be a positive random variable with a probability density function \(f(x)\) that is continuous and \(U_n=nX-[nX]\), where \([a]\) means the largest integer which is less than~\(a\). Find the limit distribution of \(U_n\) as~\(n\) goes to~\(\infty\).
\def\ProblemHeadingText{Problem 4.}
\item
Let \(X_1,X_2,\ldots,X_n\) be i.i.d. random variables with probability density function \(f(x)\), and \(M=\max(X_1,X_2,\ldots,X_n)\). Find the conditional expectation \(\mathbb E[X_1\mid M]\).
\def\ProblemHeadingText{Problem 5.}
\item
For a continuous function \(f(x)\) on \([0,1]\), define
\[
P_n(x)=\sum_{k=0}^n f\left(\frac{k}{n}\right)\binom{n}{k}x^k(1-x)^{n-k}.
\]
 Show that \(P_n(x)\) converges to \(f(x)\) on \([0,1]\) as~\(n\) goes to infinity.
\def\ProblemHeadingText{Problem 6.}
\item
Let \(S_n\) be a one-dimensional simple random walk, that is, \(S_n=X_1+\cdots+X_n\), where \(X_1,\ldots,X_n\) are i.i.d. random variables with \(\mathbb P(X_1=+1)=\mathbb P(X_1=-1)=1/2\), and \(S_0=0\). and let
\[
R_n=1+\max_{0\leq k\leq n}S_k-\min_{0\leq k\leq n}S_k
\]
 be the number of points visited by time~\(n\). Prove that \(R_n/\sqrt n\) converges weakly to some limit (in distribution).
\def\ProblemHeadingText{Problem 7.}
\item
Let \(B_t\) be standard Brownian motion with \(B_0=0\) and \(\tau=\inf\{t:B_t=a+bt\}\), where \(a>0\) and~\(b\) are two constants. Prove that
\begin{examsubparts}
\item[\textnormal{a).}]
\[
\mathbb E_0[e^{-\lambda\tau}]=\exp\{-ab-a\sqrt{b^2+2\lambda}\}
\]
\item[\textnormal{b).}]
\[
\mathbb P_0(\tau<\infty)=\exp\{-2ab\}.
\]
\end{examsubparts}
\def\ProblemHeadingText{Problem 8.}
\item
Let \(\{X_n,n=0,1,2,3,\ldots\}\) be a Markov chain on a countable irreducible state space~\(S\). Let~\(\phi\) be a nonnegative function with
\[
\mathbb E[\phi(X_1)\mid X_0=x]\leq \phi(x)
\]
 for all but a finite number of~\(x\). Prove that if \(\{x\mid \phi(x)\leq M\}\) is finite for any \(M<\infty\), then the Markov chain \(\{X_n,n=0,1,2,3,\ldots\}\) is recurrent.
\end{examproblems}
\end{document}
